% 论文的中文摘要
\begin{abstract}

随着互联网广告业务的快速发展，广告产品形态、业务规则及运营策略持续演进，围绕投放、产品说明、业务规范与运营实践等沉淀了规模庞大、类型多样的专业内容资产，这类内容在业务培训、运营决策与客户支持等场景中具有重要价值。随着业务线与内容规模持续增长，现有内容平台逐渐暴露出内容分散、跨场景复用困难、维护成本高以及质量难以评估与治理等问题，已难以满足广告业务场景的实际需求。

针对上述问题，本文以字节跳动广告业务中的实际工程项目为背景，设计并实现了一套面向广告业务的内容资产管理平台。平台以内容全生命周期为主线，对内容建模、组织、应用与治理进行系统化设计，旨在提升内容资产的统一管理能力、应用效果与质量可控性。

在内容管理方面，平台提出了一种基于空间与应用分层的内容组织模型，实现对不同广告业务线及其应用场景的隔离与统一管理，并支持文章、问答、课程和案例等多种内容形态的规范化管理。结合类目树、标签体系、权限控制与审批流程，平台降低了内容维护过程中的人工协调成本，提升了内容资产的复用性与一致性。在内容应用方面，平台在支持传统页面化内容展示的基础上，引入基于检索增强生成（Retrieval-Augmented Generation，RAG）的智能问答能力，通过内容切片、向量索引、多通道召回与重排等技术手段，为上游智能应用提供高相关性的内容支持，拓展了内容资产的使用方式。在内容治理方面，针对内容规模扩大带来的质量问题，本文构建了一种围绕重复度、歧义度与覆盖度三类问题的内容治理机制，结合向量检索与大语言模型实现对内容质量的自动检测，并通过线上化治理任务支持问题处理、复检与关闭。此外，平台通过数据洞察模块对内容使用效果和治理结果进行统计分析，为内容优化与治理策略调整提供数据支持。

应用与测试结果表明，本文所设计和实现的平台能够支撑内容管理、检索问答与治理任务处理等核心流程，并在内容定位、发布维护和治理处理等环节改善了运营效率，对广告业务场景下内容资产管理平台的建设具有一定的工程实践价值和参考意义。

% 中文关键词。关键词之间用中文全角分号隔开，末尾无标点符号。
\keywords{广告业务；内容资产管理；统一内容模型；检索增强生成；内容治理闭环}
\end{abstract}

% 论文的英文摘要
\begin{englishabstract}
With the rapid growth of Internet advertising business, advertising product forms, business rules, and operational strategies continue to evolve. A large volume of professional content has been accumulated around ad delivery, product documentation, business specifications, and operational practices, and such content plays an important role in scenarios such as business training, operational decision-making, and customer support. As business lines and content scale continue to grow, existing content platforms increasingly suffer from fragmented content sources, limited cross-scenario reuse, high maintenance costs, and insufficient capability for quality evaluation and governance.

To address these challenges, this thesis, based on a real engineering project in ByteDance's advertising business, designs and implements a content asset management platform for advertising business. Centered on the content lifecycle, the platform adopts a systematic design for content modeling, organization, application, and governance, aiming to improve the unified management capability, application effectiveness, and quality controllability of content assets.

For content consumption, the platform provides both page-based browsing and search, and introduces retrieval-augmented generation (RAG) for intelligent question answering through content slicing, vector indexing, multi-channel retrieval, and reranking, enabling upstream intelligent applications to obtain highly relevant content support and broadening the ways in which content assets can be utilized. For content quality assurance, the platform establishes an automated governance mechanism targeting duplication, ambiguity, and coverage-related issues, combining vector retrieval with large language models and operationalizing the results as online governance tasks for issue handling, rechecking, and closure. In addition, a data insight module aggregates usage and governance metrics to support iterative optimization of content and governance strategies.

The application and testing results show that the platform supports core processes such as content management, search-based question answering, and governance task handling. It also improves operational efficiency in content locating, publishing, maintenance, and governance processing. This work provides useful engineering insights and a practical reference for building content asset management platforms tailored to advertising business scenarios.

% 英文关键词。关键词之间用英文半角逗号隔开，末尾无符号。
\englishkeywords{advertising business, content asset management, unified content model, retrieval-augmented generation, closed-loop content governance}
\end{englishabstract}
